\section{问题二：变物性热湿迁移模型的建立与求解}

\subsection{变物性热湿迁移模型}

\subsubsection{非线性模型及耦合关系}

问题二保持 $R=0.02$~m、$h_T=25~\mathrm{W/(m^2K)}$ 和 $h_m=8\times10^{-7}~\mathrm{m/s}$ 不变，将\cref{eq:prop23}代入\cref{eq:heat,eq:mass}。此时温度方程的储存项和扩散项都依赖 $C$，水分方程的扩散系数同时依赖 $T$ 和 $C$。若写成算子形式，可表示为
\begin{equation}
\mathcal H[T;C]=0,\qquad \mathcal M[C;T]=0,
\label{eq:q2operator}
\end{equation}
其中两个算子通过状态相关系数双向联系。该耦合不要求在一个大矩阵中同时求解全部未知量：在每次 Picard 迭代内，先固定 $C$ 解 $\mathcal H$，再用更新后的 $T$ 解 $\mathcal M$，即可逐步逼近同一时间层的自洽状态。

与问题一相比，问题二的初始状态物性并不等于附件二常数。将 $T=28~{}^\circ\mathrm C,C=2.55~\mathrm{kg/kg}$ 代入附件三关系，得到初始密度 $976.4~\mathrm{kg/m^3}$、比热容 $3415.3~\mathrm{J/(kgK)}$、导热系数 $0.4830~\mathrm{W/(mK)}$ 和扩散系数 $5.642\times10^{-9}~\mathrm{m^2/s}$。因此，两问在相同几何和表面对流系数下仍会呈现不同响应，不能把问题一的 1800~s 结果直接作为问题二的初段结果。

\subsubsection{物性随状态的演化}

为展示非线性系数的实际变化，将初态以及 3~h 的中心、表面状态代入\cref{eq:prop23}，结果见\cref{tab:q2prop}。随着水分降低，$\rho,c_p,k$ 均减小；这意味着单位体积储热能力降低，同时内部导热能力也有所下降。二者对升温速度具有相反影响，需要由完整能量方程共同决定。

\begin{table}[!htbp]
\centering
\caption{问题二代表状态下的物性参数}\label{tab:q2prop}
\begin{tabular}{lcccc}
\toprule
状态 & $\rho$/（kg/m$^3$） & $c_p$/[J/(kg K)] & $k$/[W/(m K)] & $D$/（m$^2$/s）\\
\midrule
初始全域&976.4&3415.3&0.4830&$5.642\times10^{-9}$\\
3 h 中心&876.1&3196.9&0.4526&$1.239\times10^{-8}$\\
3 h 表面&779.0&2823.5&0.4008&$1.027\times10^{-8}$\\
\bottomrule
\end{tabular}
\end{table}

扩散系数同时受温度和水分影响。升温使 Arrhenius 因子 $\exp[-3850/(T+273.15)]$ 增大，促进水分扩散；干燥使 $\exp(-0.45/C)$ 减小，抑制扩散。3~h 时两种效应叠加后，中心和表面 $D$ 都高于初值，但表面尽管温度略高，其 $D$ 反而小于中心值。这说明低含水表层已开始形成扩散阻力，水分迁移并非随温度升高而无限加速。

\subsection{模型的求解}

\subsubsection{计算网格与非线性系数更新}

沿用问题一的 161 个表面加密节点和 1~s 时间步，但重新从 $T_i^0=28~{}^\circ\mathrm C$、$C_i^0=2.55~\mathrm{kg/kg}$ 出发，采用附件三物性计算，不能以问题一的预热结果代替初段解。前 3~h 用于本问温湿场分析，随后沿同一轨迹继续积分，为问题三确定达标时间。

对每个时间层，以 $T^{(0)}=T^n,C^{(0)}=C^n$ 为迭代初值。第 $m$ 次迭代先由 $C^{(m)}$ 更新 $\rho,c_p,k$，再以 $s_i^{(m)}=\rho_i^{(m)}c_{p,i}^{(m)}$ 和 $k_i^{(m)}$ 组装温度方程。界面系数采用相邻节点算术平均，中心零通量与表面对流条件按问题一的方法进入矩阵。

\subsubsection{温湿场交替求解}

将\cref{eq:q1implicitdetail}分别用于温度、水分，记对应三对角矩阵为 $A_T,A_C$。一次交替迭代可写为
\begin{align*}
A_T(C^{(m)})T^{(m+1)}
&=\operatorname{diag}\!\left(\frac{\rho^{(m)}c_p^{(m)}V}{\Delta t}\right)T^n+b_T,\\
A_C(T^{(m+1)},C^{(m)})C^{(m+1)}
&=\operatorname{diag}\!\left(\frac{V}{\Delta t}\right)C^n+b_C.
\end{align*}
其中 $b_T,b_C$ 为表面环境交换项。水分方程使用新温度与当前水分估计更新扩散系数，温度相关指数统一采用开尔文温度。两式右端的历史状态分别保持为 $T^n,C^n$，只在迭代收敛后推进时间。

具体步骤如下：


\textbf{Step 1 更新边界：}在 $t_{n+1}$ 读取 PCHIP 环境值；超过 14400~s 后采用 $50~{}^\circ\mathrm C$、$0.05~\mathrm{kg/kg}$ 的长期边界。

\textbf{Step 2 求解温度：}按当前水分估计更新储热系数和导热系数，求解温度三对角系统。

\textbf{Step 3 求解水分：}以新温度和当前水分估计更新扩散系数，求解水分三对角系统。

\textbf{Step 4 判断收敛：}采用问题一的温度 $10^{-8}$、水分 $10^{-10}$ 双重无穷范数容差；未满足时返回物性更新，满足后接受该时间层。



\subsubsection{步长选取与计算检查}

全隐式处理避免显式扩散格式对最小网格宽度的严格步长限制，但步长仍影响时间精度。采用 1~s 可解析附件一 60~s 采样间隔内的边界变化；长期达标时间的 2~s 与 1~s 对照相对差为 $1.52\times10^{-5}$，网格检验见后文。每步检查状态有限性、$0\le C\le2.55$ 和径向单调性，并将控制体累积变化与表面通量核对。实际最多 5 次 Picard 迭代即收敛，未触及 30 次上限。

\subsection{结果的分析和验证}

\subsubsection{温度结果及其空间演化}

典型时刻温度见\cref{tab:q2t}，空间曲线见\cref{fig:q2}(a)。0.5~h 时中心和表面分别为 $32.1911~{}^\circ\mathrm C$、$35.4141~{}^\circ\mathrm C$，仍有 $3.2230~{}^\circ\mathrm C$ 的明显差异；1.5~h 后差异降为 $1.2928~{}^\circ\mathrm C$；到 3~h 仅剩 $0.1173~{}^\circ\mathrm C$。这表明温度先经历表面快速升温和内部滞后，随后在环境逐渐稳定时趋于径向均匀。

\begin{table}[!htbp]
\centering
\caption{问题二温度分布（$^\circ$C）}\label{tab:q2t}
\begin{tabular}{cccccc}
\toprule
$t$/h&0 cm&0.5 cm&1.0 cm&1.5 cm&2.0 cm\\
\midrule
0.5000&32.1911&32.3840&32.9674&33.9620&35.4141\\
1&40.3822&40.5543&41.0609&41.8784&43.0001\\
1.5000&45.8464&45.9345&46.1927&46.6049&47.1393\\
2&48.4498&48.4877&48.5980&48.7732&49.0030\\
2.5000&49.4667&49.4790&49.5133&49.5645&49.6617\\
3&49.8493&49.8552&49.8746&49.9100&49.9666\\
\bottomrule
\end{tabular}
\end{table}

\subsubsection{水分结果及控制机制}

水分结果见\cref{tab:q2c}和\cref{fig:q2}(b)。0.5~h 时中心仍为 $2.5499~\mathrm{kg/kg}$，表面已降至 $1.6487~\mathrm{kg/kg}$；3~h 时中心和表面分别降至 $1.7662$ 和 $1.0081~\mathrm{kg/kg}$。中心—表面差由 0.5~h 的 0.9012 降至 3~h 的 $0.7581~\mathrm{kg/kg}$，说明水分扰动已经传播到中心，但内部尚未达到近似均匀状态。

\begin{table}[!htbp]
\centering
\caption{问题二水分浓度分布（kg/kg）}\label{tab:q2c}
\begin{tabular}{cccccc}
\toprule
$t$/h&0 cm&0.5 cm&1.0 cm&1.5 cm&2.0 cm\\
\midrule
0.5000&2.5499&2.5489&2.5255&2.3257&1.6487\\
1&2.5257&2.4947&2.3578&2.0230&1.4712\\
1.5000&2.3860&2.3256&2.1344&1.8020&1.3476\\
2&2.1708&2.1084&1.9236&1.6260&1.2311\\
2.5000&1.9567&1.9006&1.7360&1.4720&1.1167\\
3&1.7662&1.7165&1.5702&1.3333&1.0081\\
\bottomrule
\end{tabular}
\end{table}

3~h 时温度已接近均匀，而中心水分仍为表面的约 1.75 倍，说明升温完成不能作为烘干结束的依据。

\begin{figure}[!htbp]
\centering
\includegraphics[width=0.92\textwidth]{q2_coupled_profiles.pdf}
\caption{问题二变物性温度与水分剖面}\label{fig:q2}
\end{figure}

\subsubsection{水分迁移与扩散系数的时空证据}

\cref{fig:q2surface}展示前 3~h 水分时空分布，失水影响由表面向中心传播。

\begin{figure}[!htbp]
\centering
\includegraphics[width=0.96\textwidth]{q2_moisture_surface.pdf}
\caption{问题二前 3 h 水分浓度的三维时空曲面}\label{fig:q2surface}
\end{figure}

\cref{fig:q2diffusivity}展示 $D(T,C)$ 的时空变化：升温增强扩散，但表层失水使其扩散系数低于内部。

\begin{figure}[!htbp]
\centering
\includegraphics[width=0.96\textwidth]{q2_diffusivity_map.pdf}
\caption{问题二状态相关扩散系数的时空云图及等值线}\label{fig:q2diffusivity}
\end{figure}

\subsubsection{共享轨迹一致性检验}

问题二、三采用相同物性、网格和边界连续计算。第二问逐秒温湿场按 60~s 抽取后，与第三问全部共同时间层逐元素一致，最大差为零；长期网格与时间步检验见问题三。这一检查确认短期温湿响应与长期达标预测采用相同数值轨迹。

\subsubsection{本问结论与分析}

各时间层最多 5 次 Picard 迭代即收敛，温湿剖面保持平滑。3~h 时中心—表面温差为 $0.1173~{}^\circ\mathrm C$、水分差为 $0.7581~\mathrm{kg/kg}$，后续干燥由内部水分迁移控制。
